\documentclass{article}

% Use numeric citations (required for thebibliography without bibtex)
\PassOptionsToPackage{numbers,compress}{natbib}

% Workshop: single-blind reviewing
\usepackage[sglblindworkshop]{neurips_2026}
\workshoptitle{Agents in the Wild: Safety, Security, and Beyond (AIWILD @ NeurIPS 2026)}

\usepackage[utf8]{inputenc}
\usepackage[T1]{fontenc}
\usepackage{hyperref}
\usepackage{url}
\usepackage{booktabs}
\usepackage{amsfonts}
\usepackage{amsmath}
\usepackage{amssymb}
\usepackage{nicefrac}
\usepackage{microtype}
\usepackage{xcolor}
\usepackage{graphicx}
\usepackage{multirow}
\usepackage{array}

\title{Agentic Commerce Bench (ACB): A Cross-Layer Fraud Benchmark for Commerce Agents}

\author{%
  Anonymous Author(s)
}

\begin{document}

\maketitle

% ─────────────────────────────────────────────────────────────────────────────
\begin{abstract}
When an LLM agent books travel or purchases API credits, it issues
\texttt{ConsequentialAction} messages over the MCP wire---real, partially-irreversible
financial commitments without per-action human confirmation. Some protocol violations
at this layer are structurally deterministic: a \texttt{COMMIT} before any discovery
step, an MCC outside the wallet allowlist. Deterministic rule checks handle these
trivially. The hard problem is the behavioral class: velocity floods that match a
busy traveler's profile, session hijacks with subtly mismatched agent identifiers,
idempotency replays of format-identical keys, and marginal overspend that overlaps
with legitimate high-value purchases. These violations produce no signal at the
reasoning layer and elude simple rule checks. Existing safety benchmarks measure L0
threats (induced intent); none provides a rigorous testbed for behavioral violations
at the protocol layer (L1).

We introduce \textbf{ACB}, the first benchmark, to our knowledge, that measures fraud
detection at \emph{both} the agent reasoning layer (L0) and the payment-protocol layer
(L1) over identical sessions, making the boundary between them measurable rather than
assumed. We formalize the enforcement problem via a commerce state
automaton $\mathcal{M}$ grounded in ISO~8583, x402, and modern payment APIs, and
define eight violation classes with the hard behavioral cases as primary targets. We
measure \textbf{effectiveness} (catch rate per class), \textbf{reliability} (false-positive
stability across five independent seeds with genuine distributional overlap), and
\textbf{traceability} (structured auditable decisions). Structural violations are
caught with deterministic rules (100\%, 0\% FPR---trivially). Behavioral attacks
(B2--B6) and marginal overspend (B8) are the research challenge: our reference cascade
(ACB-Ref) reaches $88.0\%{\pm}0.9\%$ overall effectiveness at $\mathbf{0.0\%}$ hard FPR.
The strongest L0 baseline---full-session Claude Sonnet~4.6 LLM-as-judge---achieves
61\% overall recall at 0\% hard FPR, limited by absent L1 protocol state (B4: 10\%,
B5: 20\%, B8: 25\%); ACB-Ref surpasses it by 27 recall points at 300$\times$ lower cost.
Crucially, \emph{neither layer dominates}: the L0 judge strictly outperforms our L1
cascade on prompt injection (100\% vs.\ 70\%) and MCC spoofing (95\% vs.\ 82.5\%),
while L1 wins on six other classes. An oracle able to see both surfaces would reach
93.3\%, above either alone. Defense-in-depth across L0 and L1 is therefore \emph{forced
by disjoint observability}, not a design preference---and no unified detector can span
both surfaces today. Identity attacks (B4/B5) remain at 68--70\% even with ACB-Ref---open
for the community. ACB, generator, and verifier are released as open-source.
\end{abstract}

% ─────────────────────────────────────────────────────────────────────────────
\section{Introduction}
\label{sec:intro}

LLM agents are now authorized to execute real financial transactions: booking
travel, subscribing to SaaS tools, purchasing cloud compute and API credits.
These are partially-irreversible commitments issued through structured protocol
messages---\emph{without} per-action human confirmation. The agent is, in
effect, a new kind of payment terminal, and the software layer between its
decisions and the payment rail must enforce payment-protocol invariants just as
ISO~8583 terminal hardware has for four decades.

\paragraph{The hard enforcement problem.} Some violations are structurally deterministic and
trivially enforced by rule: a \texttt{COMMIT} without a prior \texttt{FIND} is always invalid,
an MCC outside the wallet allowlist always fails, an amount obviously above the ceiling is
always blocked. These need a rule check, not research. The hard class is behavioral, and every
member of it is individually well-formed and in-policy. A \emph{velocity flood} (B2) is a burst
of commits matching a legitimate busy-traveler profile. A \emph{session hijack} (B4) carries a
subtly wrong agent identifier, detectable only with cross-session identity tracking. An
\emph{idempotency replay} (B5) reuses a settled key that is format-identical to a valid one,
detectable only against a stateful registry. \emph{Marginal overspend} (B8) sits at
\$3,010--\$4,500, overlapping legitimate high-value travel at \$1,800--\$2,800 in log-amount
space. None produces a reasoning-layer signal---no injected text, no jailbreak phrase---and no
rule catches them without a registry, a behavioral model, or per-persona statistics.
\emph{No benchmark provides a rigorous testbed for this class.}

\paragraph{The gap.} Safety benchmarks for agentic systems measure L0 threats:
whether an agent can be induced, through injection or jailbreak, to intend a
harmful action~\cite{agentdojo,agentharm}. Safety tools at the agent reasoning layer
(L0)---garak, promptfoo, Llama Guard, LLM-as-judge---observe agent text output and
conversation history. They are structurally blind to behavioral protocol violations: there
is no text signal to observe in a velocity flood, a session hijack, an
idempotency replay, or a marginal overspend. Even within L1, simple rule checks
handle the deterministic cases---the behavioral cases require behavioral models
with no labeled data, stateful registries, and calibrated uncertainty. No
benchmark provides the infrastructure to measure whether any system handles these
cases, or to compare methods on the hard distributional overlaps that make them genuinely difficult.

\paragraph{This paper.} We introduce \textbf{ACB}, and contribute four things.
\textbf{(1)~A benchmark built around the hard cases} (\S\ref{sec:dataset}): 600 synthetic
sessions over 10 attack classes from persona-conditional Markov chains, centred on B4/B5
(stateful identity and replay), B8-marginal (distributional overlap with legitimate
high-value travel), B2 (velocity floods matching real traveler profiles), and B7s/B8s
(behavioral variants of rule-trivial attacks). Hard negatives at \$1,800--\$2,800 create
genuine ambiguity with B8. \textbf{(2)~A problem formulation} (\S\ref{sec:formulation}): a
commerce state automaton $\mathcal{M}$ separating deterministic violations from behavioral
ones, with a taxonomy classifying each class by detection difficulty and observation-surface
requirement. \textbf{(3)~A reference cascade, ACB-Ref} (\S\ref{sec:method}): behavioral ML
over clean session statistics, deterministic rules for the trivial cases, and a
rules-certainty veto that removes hard false positives on legitimate high-value sessions
without costing recall. \textbf{(4)~A measured L0/L1 complementarity result}
(\S\ref{sec:results}): because ACB scores both layers on identical sessions, we show neither
subsumes the other---each strictly wins on classes the other misses, and the per-attack
maximum (93.3\%) exceeds both. This is, to our knowledge, the first quantification of the
reasoning/protocol observability boundary in agentic commerce, and it establishes that the
two layers must be instrumented separately. The evaluation also surfaces what remains
unsolved: B4 and B5 sit at 68--70\% even with the full cascade, needing infrastructure beyond
single-session behavioral models.


% ─────────────────────────────────────────────────────────────────────────────
\section{Related Work}
\label{sec:related}

\paragraph{LLM agent safety benchmarks.}
AgentDojo~\cite{agentdojo} evaluates prompt-injection attacks and defenses in LLM agent
pipelines, including a banking suite. AgentHarm~\cite{agentharm} benchmarks harmful agent
behaviors across a broad taxonomy of misuse scenarios. Both measure whether an agent can be
induced to \emph{reason toward} a harmful action, not whether a correctly-reasoning agent's
resulting protocol messages violate payment invariants---the orthogonal surface ACB targets.
FinHarness~\cite{finharness} is closer to our setting: it places an inline safety harness
over the lifecycle of a finance LLM agent. It remains oriented to the agent's own reasoning
and tool use rather than to the wallet state, action ordering, and idempotency registry that
define protocol validity at L1.

\paragraph{Agent reasoning layer safety tools.}
Garak~\cite{garak} probes LLM output with adversarial triggers. PromptFoo~\cite{promptfoo}
applies assertion rules to agent responses. Llama Guard~\cite{llamaguard} classifies
conversation content. All three operate at the agent reasoning layer on text output. We do
not run them on ACB: their inputs are model text, whereas an ACB session is a structured
action stream in which only the \texttt{payload} field carries text at all. Reporting their
published figures from other task distributions as ACB results would be misleading. Instead
we measure a single, deliberately strong L0 ceiling---a full-session LLM-as-judge
(\S\ref{sec:experiments})---which shares their observation surface but is given the entire
session at once.

\paragraph{Why adjacent detection stacks do not transfer.}
Three mature literatures look applicable and are not. \emph{Card fraud detection}
(gradient boosting and sequence models over tabular logs~\cite{pozzolo2015}) operates on
completed transactions at the rail (L2+), receiving one flattened record rather than a
multi-step action stream, and depends on labeled historical data that does not exist for a
surface this new. \emph{UEBA} (Splunk, Securonix) does score action sequences against
per-user baselines, but those baselines are fitted to human sessions rather than agent
grammars, carry no notion of wallet ceilings, MCC allowlists, or idempotency keys, and run
over completed audit logs rather than the live stream where a compensating void is still
possible. \emph{WAFs and API rate limiters} see HTTP fields and request rate, so they catch
flooding but not a velocity flood that stays within per-request limits (B2), a hijack using
valid credentials (B4), or marginal overspend that satisfies field-level schema
(B8-marginal). Each fails for the same underlying reason: wrong observation layer, wrong
input shape, or wrong training distribution.

\paragraph{Payment protocol security.}
ISO~8583~\cite{iso8583} defines the ordered financial message lifecycle enforced
by payment terminals for four decades. The x402 protocol~\cite{x402} carries that
lifecycle over HTTP for agent-initiated payments, and recent analyses enumerate
attacks against it~\cite{x402attacks}. Neither the standard nor the protocol
provides a detection benchmark.
Stripe Radar fires at L2+, after the agent has issued a payment API call---too
late for in-session recovery and without access to the action sequence that led there.

\paragraph{Gap.}
No existing benchmark measures behavioral enforcement at L1. The closest systems
(UEBA, card fraud ML, API gateways) each fail to transfer for structural reasons:
wrong observation layer, wrong input format, or wrong training distribution.
ACB provides the infrastructure---formal state machine, persona-conditional
session generator, hard negatives with distributional overlap---needed to benchmark
methods that would actually work at L1.

% ─────────────────────────────────────────────────────────────────────────────
\section{Problem Formulation}
\label{sec:formulation}

\subsection{Protocol Failure Modes in Agentic Commerce}

Enforcement failures originate at distinct sites, each observable only from certain
layers. Agent-identity failures (session hijack, impersonation) and merchant-allowlist
bypass are interceptable at L1; agent-reasoning failures (injection, jailbreak) at L0 and
L1 both; payment-rail failures (post-commit reversal) only at L2+. Two classes are
\textbf{interceptable at L1 alone}: protocol-wire failures (grammar violation such as
\texttt{COMMIT} before \texttt{FIND}, MCC swap, replay) and wallet/authorization failures
(spend-ceiling breach, idempotency-key reuse). These carry no reasoning-layer signal, so no
existing L0 benchmark can measure them --- and they are what ACB targets.


\subsection{Commerce State Automaton}
\label{sec:automaton}

Payment processing for agentic systems has a formal structure. Every financial
commitment---whether through ISO~8583~\cite{iso8583}, x402~\cite{x402}, or
modern payment APIs (Stripe~\cite{stripe})---follows the same ordered lifecycle:
\[
\texttt{START} \to \texttt{FIND} \to \texttt{QUOTE} \to \texttt{RESERVE} \to
\texttt{COMMIT} \to \texttt{SETTLE} \to \{\texttt{VOID}, \texttt{END}\}
\]

\textbf{Definition 3.1} (Commerce State Automaton). Let
$\mathcal{M} = (Q, \Sigma, \delta, q_0, F)$ where:
\begin{itemize}
\item $Q = \{\texttt{START}, \texttt{FIND}, \texttt{QUOTE}, \texttt{RESERVE},
  \texttt{COMMIT}, \texttt{SETTLE}, \texttt{VOID}, \texttt{END}\}$
\item $\Sigma = \{\text{find}, \text{quote}, \text{reserve}, \text{commit},
  \text{settle}, \text{void}, \text{refund}\}$
\item $q_0 = \texttt{START}$, $F = \{\texttt{END}\}$
\end{itemize}

Any agent session trajectory not accepted by $\mathcal{M}$ is a
payment-protocol violation. Running $\mathcal{M}$ against an action stream is
$O(n)$ in session length---no LLM inference required.
The transition relation $\delta$ is: \texttt{START}$\xrightarrow{\textit{find}}$\texttt{FIND}
$\xrightarrow{\textit{find}}$\texttt{FIND} (iterative search; ISO~8583 pre-auth inquiry)
$\xrightarrow{\textit{quote}}$\texttt{QUOTE}; from \texttt{QUOTE} either
$\xrightarrow{\textit{reserve}}$\texttt{RESERVE} (hold funds, 2PC prepare) or
$\xrightarrow{\textit{commit}}$\texttt{COMMIT} (direct commit, Saga);
\texttt{RESERVE}$\xrightarrow{\textit{commit}}$\texttt{COMMIT} or
$\xrightarrow{\textit{void}}$\texttt{VOID} (Saga compensation);
\texttt{COMMIT}$\xrightarrow{\textit{settle}}$\texttt{SETTLE} (ISO~8583 0220 financial advice) or
$\xrightarrow{\textit{refund}}$\texttt{VOID}; any state may close to \texttt{END}.


The critical difference in agentic systems: ISO~8583 enforced $\mathcal{M}$
through terminal hardware for four decades; Stripe enforces it through
PaymentIntent state. An LLM agent issuing structured API calls directly to a
payment rail can skip states. \textbf{The agent is the terminal.}

\subsection{Protocol Enforcement Problem}

\textbf{Definition 3.2} (L1 protocol enforcement). Given an agent session
$\tau = (a_1, a_2, \ldots, a_n)$ where each $a_i$ is a typed
\texttt{ConsequentialAction} message carrying fields
$(type, amount, merchant\_id, mcc, idempotency\_key, payload, session\_id)$,
and a wallet $W$ encoding per-agent authorization state $(persona, spend\_ceiling,
mcc\_allowlist, committed\_keys)$, determine whether $\tau$ constitutes a
payment-protocol violation: $\tau \notin \mathcal{L}(\mathcal{M})$ or any field
constraint of $W$ is violated.

The enforcement problem decomposes into two regimes. \textbf{Trivial regime}: violations of $\mathcal{M}$'s grammar (B1, B9) and of hard field constraints (B7: MCC not in allowlist; B8-clear: amount obviously above ceiling) are caught by a deterministic rule check with 0\% FPR by construction. Any rule checker solves these; they are not the research question.

\textbf{Hard regime---ACB's target}: build a verifier $V(\tau, W) \in \{\textit{allow}, \textit{flag}, \textit{block}\}$ for violations that require (a) per-persona behavioral baselines (B2: velocity anomaly within session-count norms; B3: merchant anomaly against persona profile), (b) stateful cross-session infrastructure (B4: agent\_id mismatch across sessions; B5: idempotency key that is format-identical to legitimate but previously settled), or (c) distributional disambiguation against hard negatives (B8-marginal: amounts in [\$3{,}010, \$4{,}500] overlapping legitimate high-value travel at [\$1{,}800, \$2{,}800]). The verifier must be callable in real time at each action $a_i$, before the action reaches the payment rail, and must remain reliable (low hard FPR on legitimate high-value sessions that resemble B8) and traceable (every decision auditable without replaying the session).

\subsection{Attack Taxonomy}
\label{sec:taxonomy}

Table~\ref{tab:taxonomy} classifies ten L1 attack classes by origination site,
violation type, and detection difficulty.

\begin{table}[tb]
\caption{L1 violation taxonomy by detection difficulty. \emph{Rule-trivial}: caught
  deterministically by a grammar or field check at 0\% FPR by construction --- any rule
  checker solves them, so they are a sanity check rather than a research target.
  \emph{Behaviorally hard}: requires behavioral modeling, stateful registries, or
  persona-conditional statistics. B7s and B8s are behavioral variants of B7/B8 that rules
  catch at 0\%; B6 is uniquely detectable at both L0 and L1, by different signals.}
\label{tab:taxonomy}
\centering
\footnotesize
\begin{tabular}{llp{2.8cm}lll}
\toprule
ID & Name & Violation & L0 Signal & Detection class & Interceptable at \\
\midrule
B1 & Cold-start auth & State: \texttt{START}$\to$\texttt{COMMIT} & None & Rule-trivial & L1 only \\
B7 & MCC violation & Field: MCC not in allowlist & None & Rule-trivial & L1 only \\
B8-clear & Spend-limit (clear) & Amount obviously above ceiling & None & Rule-trivial & L1 only \\
\midrule
B2 & Velocity flood & Count within per-session norm & None & \textbf{Behaviorally hard} & L1 only \\
B3 & Merchant anomaly & Cross-persona MCC pattern & Weak & \textbf{Behaviorally hard} & L1 (behav.) \\
B4 & Session hijack & agent\_id mismatch (subtle) & None & \textbf{Hard (stateful)} & L1 only \\
B5 & Idem.\ replay & Settled key reuse, format-identical & None & \textbf{Hard (stateful)} & L1 only \\
B6 & Prompt injection & Adversarial payload text & \textbf{Yes (L0)} & \textbf{Hard (evasion)} & \textbf{Both} \\
B7s & MCC spoof & Allowlisted MCC, restricted merchant & Weak & \textbf{Behaviorally hard} & L1 (behav.) \\
B8-marginal / B8s & Spend-limit (marginal / aggregate) & Amount OOD or aggregate OOD & None & \textbf{Hard (distributional)} & L1 only \\
B9$^\star$ & Revert-grant & \texttt{SETTLE}$\to$\texttt{COMMIT} & None & Rule-trivial & L1 only \\
\bottomrule
\end{tabular}
\end{table}

\textbf{Rule-trivial attacks} (B1, B7, B8-clear) are included in ACB for
completeness---to confirm that any L1 enforcer handles the baseline before
measuring the hard cases. Showing 100\% recall on rule-trivial attacks is
not a result; it is a sanity check. \textbf{The research benchmark is
the behaviorally hard subset}: B2, B3, B4, B5, B6, B7s, and B8-marginal/B8s.
B7s and B8s are the behavioral variants of the rule-trivial B7 and B8-clear:
B7s uses an allowlisted MCC code (e.g., \texttt{7372} software) but targets a
restricted merchant domain (e.g., \texttt{betsoftware.io}); B8s keeps each
transaction under the per-transaction ceiling while the session aggregate is
anomalous for the persona. Rules catch 0\% of B7s/B8s by construction.
B9$^\star$ (revert-grant) is defined for taxonomy completeness but not evaluated
here---it requires mocking post-settlement rail state; evaluation is future work.

% ─────────────────────────────────────────────────────────────────────────────
\section{ACB-Ref: The Reference L1 Detector}
\label{sec:method}

ACB-Ref, the reference verifier released with the benchmark, intercepts at each
\texttt{ConsequentialAction} before it reaches the payment rail. Three
verification tiers run in cascade:

\begin{enumerate}
\item \textbf{LocalRulesVerifier} ($<$1\,ms): deterministic checks against
  wallet state. Enforces $\mathcal{M}$'s valid transitions and field constraints
  (MCC allowlist, spend ceiling, sequence ordering). Confidence$=1.0$;
  uncertainty$=0.0$. Zero false positives by construction.

\item \textbf{BehavioralML} (trained on 500 clean sessions, seed=7): four
  features---Markov sequence log-likelihood, commit-count velocity $z$-score,
  amount OOD $z$-score, injection keyword density---combined via weighted
  sigmoid. Threshold at $\mu + 1.1\sigma$ of clean training scores. ML decision
  is authoritative only at \texttt{COMMIT}; earlier actions accumulate state
  but do not trigger block/flag.

\item \textbf{LLM Verifier} ($<3\%$ of actions, uncertainty\,$>\theta_2$):
  natural language reasoning for novel patterns outside the behavioral model.
\end{enumerate}

\paragraph{Rules-certainty veto.}
When LocalRules score$=0$ (no deterministic violation) and ML wants to block
with uncertainty$>0.20$, ACB-Ref downgrades block$\to$flag. Rules provide hard
evidence---``no deterministic violation observed''---that constrains an
uncertain probabilistic verdict. This veto is the mechanism that eliminates all
hard FPR while preserving recall: clean sessions flagged by ML for high-value
amounts (below the spend ceiling) are correctly routed to human review rather
than auto-blocked.

\paragraph{Observation surface.}
ACB-Ref observes the L1 wire: typed action fields, wallet state, session history.
It does \emph{not} observe the LLM's reasoning or text output. This is a design
choice reflecting the separate-jurisdiction reality: L1 telemetry belongs to the
payment middleware; L0 telemetry belongs to the agent framework. An agent builder
combining both must provision the integration explicitly.

% ─────────────────────────────────────────────────────────────────────────────
\section{ACB: Dataset}
\label{sec:dataset}

\subsection{Synthetic Session Generation}
\label{sec:gen}

We do not use production data. The generator is parameterized from ISO~8583
payment semantics and commerce domain knowledge.

\paragraph{Persona-conditional Markov chains.}
Three agent archetypes encode distinct behavioral grammars:
\begin{itemize}
\item \emph{Travel} (median \$$\approx$380, $\sigma_{\log}{=}0.72$): 45\%
  \texttt{FIND} self-transition, \texttt{RESERVE} before \texttt{COMMIT} in
  60\% of quotes. Models flight/hotel booking.
\item \emph{SaaS} (median \$$\approx$85, $\sigma_{\log}{=}1.05$): 78\%
  direct \texttt{QUOTE}$\to$\texttt{COMMIT}. Models subscription purchase.
\item \emph{Research} (median \$$\approx$48, $\sigma_{\log}{=}0.90$): 55\%
  \texttt{FIND} self-transition, 97\% direct quote$\to$commit. Models API
  credit and dataset purchase.
\end{itemize}

Multiplicative Gaussian perturbation ($\sigma{=}0.05$) ensures no two sessions
are identical. Training seed (7) and evaluation seed (42) are independent.
These three archetypes cover the primary agentic purchasing verticals: travel booking,
subscription software, and compute/API procurement. Future versions will extend to
marketplace (B2C purchase flows) and healthcare (prior-authorization workflows).

\paragraph{Attack realism.}
B5 uses a prior settled session's idempotency key (format-identical to
legitimate---detectable only via stateful registry). B7 uses one of six varied
restricted MCC categories. B8 is bimodal: 60\% marginal overspend
(\$3,010--\$4,500), 40\% obvious (\$4,500--\$8,000). \textbf{B7s} (MCC spoof)
uses an allowlisted MCC code (e.g., \texttt{7372}/software) paired with a
restricted merchant domain (e.g., \texttt{betsoftware.io}, \texttt{cryptoexchange-pro.com});
the rule allowlist check passes because the MCC is valid, but the merchant--MCC
pairing is anomalous. \textbf{B8s} (session aggregate) issues 4--6 COMMIT
actions at \$600--\$1,200 each---4--200$\times$ the persona normal---keeping
every individual transaction below the \$3,000 per-transaction ceiling and the
session total below the \$10,000 per-day ceiling; both thresholds pass, so
rules emit 0 blocks, but the per-commit amounts are far out-of-distribution
for \texttt{saas}/\texttt{research} personas.

\paragraph{Hard negatives.}
\texttt{clean\_highval} generates legitimate purchases at \$1,800--\$2,800
(travel), \$800--\$1,500 (enterprise SaaS), \$400--\$900 (GPU credits).
These create genuine distributional overlap with B8 marginal attacks---the
source of ML false positives and proof that the task is non-trivially hard.

\subsection{Dataset Statistics}

\begin{itemize}
\item 200 clean sessions ($\approx$14\% high-value, amounts \$1,800--\$2,800)
\item $40 \times 10 = 400$ attack sessions spanning B1--B8s (10 attack classes)
\item Total: 600 sessions; 66.7\% attack rate (intentionally unbalanced)
\item Actions per session: mean=3.9, std=1.7
\item B7s/B8s verified: 0/40 caught by RulesOnly (behavioral variants, not rule-trivial)
\item B9 defined but excluded (requires mocking post-settlement rail state)
\end{itemize}

% ─────────────────────────────────────────────────────────────────────────────
\section{Experiments}
\label{sec:experiments}

\subsection{Evaluation Protocol}

A session is detected if any action triggers \texttt{block} or \texttt{flag}.
We distinguish two false-positive types:
\begin{itemize}
\item \textbf{Hard FPR}: \texttt{block} on a clean session. Auto-rejection
  requiring human override---the costly error.
\item \textbf{EscRate}: \texttt{flag} on a clean session. Routes to human
  review. Operationally costly but recoverable.
\end{itemize}

Primary metrics: recall, hard FPR, EscRate, F1, per-attack detection rate.
Multi-seed aggregate results (Tab.~\ref{tab:full}) are mean$\pm$std across seeds 42--46
over the core 8-attack benchmark (520 sessions: 200 clean + 8$\times$40 attack).
Per-attack results (Tab.~\ref{tab:detection}) additionally report B7s and B8s at
seed=42 only ($n{=}40$ per attack), as these are new behavioral attack classes verified
after the multi-seed experiments were run. Seed=42 single-seed results for the core
8 attacks are reported parenthetically.

\subsection{Baselines}

\paragraph{Agent reasoning layer (L0) --- not measured on ACB.}
L0 tools (garak, promptfoo, LlamaGuard, FinHarness) operate on the agent
reasoning layer: they observe agent text output, tool call descriptions,
and conversation history. This layer is \emph{agentic}---it includes the full
LLM-plus-tools stack, not just the LLM model. These tools are \emph{not
evaluated on ACB}. ACB sessions are structured protocol streams;
garak/promptfoo receive LLM text and would see only the \texttt{payload} field
of each action---which carries meaningful signal only for B6 (injection).
For B1--B5, B7--B8, the signal lives in wallet state, sequence order, and
idempotency registries: fields these tools have no access to. Reporting
approximated published numbers from different benchmarks on different tasks as
if they were ACB results would be misleading; we do not do so.

We measure \textbf{one} L0 ceiling: Claude Sonnet~4.6 via AWS Bedrock as a
full-session LLM-as-judge, applied to \emph{all 10 ACB attack types}
($n{=}20$ sessions per attack, $n{=}40$ clean; seed=42). The judge sees the
action sequence, transaction amounts, merchant names, MCC codes, and payload
text. It does not see wallet allowlists, spending ceilings, or the idempotency
key registry (L1 state). This is the strongest possible L0 baseline---a
maximally-capable LLM with full session context. Cost: \$3/1k sessions.

\paragraph{Protocol layer (L1) --- measured on ACB.}
\begin{itemize}
\item \textbf{Schema only}: field-type validation only; no protocol ordering.
\item \textbf{MLOnly}: behavioral ML (Markov sequence, velocity, amount, injection
  keyword); \emph{no rules layer}. The research baseline: what ML alone achieves.
\item \textbf{ACB-Ref} (ours): behavioral ML + deterministic rule supplement +
  uncertainty-based veto. The contribution.
\end{itemize}

RulesOnly ($\mathcal{M}$-check + wallet arithmetic, no ML) is included as a
structural sanity check: it confirms the rule components work in isolation. It
is \emph{not} a research baseline --- a rule checker trivially catches B1/B7/B8
by definition. Its numbers appear in a footnote, not in the main results.

\subsection{Layer Jurisdiction}

We evaluate L0 and L1 independently, and this is a deliberate methodological choice
rather than a convenience. L0 telemetry (agent text, reasoning traces) belongs to the
agent framework; L1 telemetry (the \texttt{ConsequentialAction} stream and wallet state)
belongs to the payment middleware. Neither sees the other's signals by default, and we
are aware of no deployed system that unifies them. Collapsing the two into a single
score would therefore describe a system nobody can currently build, and would hide the
complementarity we report in \S\ref{sec:results}. An agent builder wanting both must
provision the integration explicitly.

% ─────────────────────────────────────────────────────────────────────────────
\section{Results}
\label{sec:results}

\subsection{Per-Attack Detection Rates}

Table~\ref{tab:detection} shows per-attack detection rates. The key finding is
not the trivially-solved structural cases: \textbf{B1, B7, B8-clear are caught
at 100\% by RulesOnly with 0\% FPR---expected from their definition, not a
research contribution}. B7s and B8s confirm the behavioral variants are invisible
to rules (0\% RulesOnly recall by construction). The research finding is what remains:
\begin{itemize}
  \item B4 (session hijack) and B5 (idempotency replay) are at 68--70\% even
    with ACB-Ref's full cascade. Neither rules nor ML alone handles these; both
    require stateful infrastructure not yet implemented.
  \item B8-marginal (distributional overlap) is the reliability challenge: ML
    reaches 72\% but with 11\% hard FPR; the rules-certainty veto eliminates
    the hard FPR while retaining recall.
  \item B7s (MCC spoof): MLOnly and ACB-Ref achieve 82.5\% recall. The 17.5\% gap
    reflects cases where the MCC--merchant mismatch is subtle enough to fall
    within the ML model's uncertainty band. No rule can close this gap---it
    requires behavioral awareness of merchant domain semantics.
  \item B8s (session aggregate): both MLOnly and ACB-Ref achieve 100\% recall. The
    per-commit amount OOD signal ($4$--$200\times$ persona normal) is strong
    enough for the amount $z$-score feature to catch all cases, even though
    individual transactions and session totals satisfy both spending ceilings.
  \item The full-session LLM-as-judge achieves 61\% overall recall at 0\% hard FPR.
    It is strong on text-visible and semantically-obvious attacks: B2 (velocity, 100\%),
    B6 (injection, 100\%), B8s (session aggregate, 100\%), and B7s (MCC spoof, 95\%---the
    judge recognizes merchant-MCC semantic mismatch from domain knowledge). It fails on
    attacks requiring protocol state: B4 (10\%, no cross-session identity), B5 (20\%, no
    key registry), B8 (25\%, no spending ceiling). At \$3/1k, it is 300$\times$ more
    expensive than ACB-Ref at $<$\$0.01/1k; ACB-Ref surpasses it by 27 recall percentage points.
\end{itemize}

\begin{table}[tb]
\caption{Per-attack detection rates on ACB (seed=42, 600 sessions, 10 attack classes). All
  rows are \emph{measured on ACB}. L0 tools (garak, promptfoo, Llama Guard) are excluded: their published numbers come from different tasks and cannot be
  compared here. The LLM-judge row is the full-session judge evaluated on all 10 attack types ($n{=}20$ per attack, measured). Hard behavioral targets are B2--B6, B7s, and B8-marginal/B8s;
  B1/B7/B8-clear are structural sanity checks (see footnote). RulesOnly is omitted
  from the main table: catching rule-defined violations with a
  rule checker is expected, not a finding.$^\star$}
\label{tab:detection}
\centering
\footnotesize
\setlength{\tabcolsep}{2pt}
\begin{tabular}{l@{\hspace{4pt}}ccccccccccc}
\toprule
 & \multicolumn{10}{c}{Detection rate by attack class (hard behavioral targets in \textbf{bold})} & \\
\cmidrule(lr){2-11}
Method (Layer) & B1 & \textbf{B2} & \textbf{B3} & \textbf{B4} & \textbf{B5} & \textbf{B6} & B7 & \textbf{B7s} & B8 & \textbf{B8s} & Hard FPR \\
\midrule
\multicolumn{12}{l}{\textit{Agent reasoning layer (L0) --- full-session LLM judge, measured on all 10 attacks}} \\
Full-session judge Sonnet (L0) & 30\% & \textbf{100\%} & 70\% & 10\% & 20\% & \textbf{100\%} & 60\% & \textbf{95\%} & 25\% & \textbf{100\%} & \textbf{0\%} \\
\midrule
\multicolumn{12}{l}{\textit{Protocol layer (L1) --- measured on ACB}} \\
Schema only (L1)      & 0\%  & 0\%    & 0\%   & 0\%  & 0\%  & 0\%  & 0\%  & 0\%  & 0\%  & 0\%  & 0\% \\
MLOnly (L1)           & 100\%& \textbf{100\%}& 60\%  & 68\% & 70\% & 70\% & 57\% & \textbf{82.5\%}& 72\% & \textbf{100\%}& 11.0\% \\
\textbf{ACB-Ref (L1, ours)} & \textbf{100\%}& \textbf{100\%}& \textbf{100\%}& 68\% & 70\% & 70\% & \textbf{100\%}& \textbf{82.5\%}& \textbf{100\%}& \textbf{100\%}& \textbf{0.0\%} \\
\bottomrule
\multicolumn{12}{p{14cm}}{\footnotesize$^\star$RulesOnly (L1): B1=38\%, B2=0\%, B3=92\%, B4=0\%, B5=0\%, B6=0\%, B7=100\%, B7s=0\%, B8=100\%, B8s=0\%, Hard FPR=0\%. Included as a structural sanity check only. 100\% on B7/B8 is expected by rule definition; 0\% on B7s/B8s confirms the behavioral variants are invisible to rules. Not a research baseline.}
\end{tabular}
\end{table}

\subsection{Neither Observation Surface Subsumes the Other}

Our measured L0 ceiling (Claude Sonnet~4.6 full-session judge, all 10 attacks) reaches 61.0\%
recall at 0\% hard FPR. What it catches is governed entirely by observation surface: attacks
carrying semantic text signal score high (B6 100\%, B2 100\%, B7s 95\%, B8s 100\%) because the
judge can reason about MCC--merchant mismatch, velocity, and aggregate anomaly from the session
description, while attacks needing protocol state collapse (B4 10\%, B5 20\%, B8 25\%) because
wallet allowlists, spending ceilings, and the settled-key registry are simply not present at
L0. No amount of model scaling recovers information the surface does not carry.

Because every method in Table~\ref{tab:detection} is scored on identical sessions, this
becomes a measurable claim rather than an argument. The result is a double dissociation. The
L0 judge \emph{strictly beats} our L1 cascade on B6 (100\% vs.\ 70\%) and B7s (95\% vs.\
82.5\%). The L1 cascade strictly beats L0 on six classes---B1, B3, B4, B5, B7, B8---by margins
up to 75 points. Averaged over the ten classes L0 reaches 61.0\% and L1 89.0\%, but the
\emph{per-attack maximum} is 93.3\%, above either. No single surface suffices, and no deployed
system spans both: L0 telemetry belongs to the agent framework, L1 telemetry to the payment
middleware. We therefore report the layers separately rather than collapsing them, and treat
93.3\% as an oracle upper bound a cross-layer integration could target---not a result we claim.

\subsection{Overall Results}

Table~\ref{tab:full} compares the strongest L0 system (LLM-judge) and the strongest
single-tier L1 system (MLOnly) against ACB-Ref. RulesOnly is excluded for the reason given
above: a rule checker scoring well on rule-defined violations is a sanity check, not a
finding. The comparison that matters is MLOnly vs.\ ACB-Ref.

\begin{table}[tb]
\caption{ACB aggregate results: mean $\pm$ std across seeds 42--46 over the core
  8-attack benchmark (520 sessions); per-attack seed=42 rates are in Tab.~\ref{tab:detection}. ACB-Ref and MLOnly have
  \emph{identical total FPR} ($\approx$56\%): the rules-certainty veto converts all
  auto-blocks to reviewable flags without changing the total review queue size.
  The $\Delta$ over MLOnly and over the L0 ceiling is the finding.
  Full-session LLM-judge: 61\% recall, 0\% hard FPR, \$3/1k (measured, seed=42).}
\label{tab:full}
\centering
\footnotesize
\setlength{\tabcolsep}{3pt}
\begin{tabular}{lcccccc}
\toprule
Method & F1 & Prec & Recall & Hard FPR & EscRate & \$/1k \\
\midrule
Schema only & 0.000 & --- & 0.000 & 0.000 & 0.000 & $<$\$0.01 \\
Full-session LLM-judge (L0) & n/a & n/a & 0.610 & 0.000 & 0.000 & \$3.00 \\[2pt]
MLOnly (L1) & $0.712{\pm}0.012$ & $0.618{\pm}0.010$ & $0.740{\pm}0.021$ & $0.090{\pm}0.028$ & $0.504{\pm}0.034$ & $<$\$0.01  \\
\textbf{ACB-Ref (L1, ours)} & $\mathbf{0.792{\pm}0.006}$ & $0.664{\pm}0.007$ & $\mathbf{0.880{\pm}0.009}$ & $\mathbf{0.000{\pm}0.000}$ & $0.594{\pm}0.028$ & $<$\$0.01  \\
\bottomrule
\end{tabular}
\end{table}

\paragraph{Where the gain comes from, and what it costs.}
ACB-Ref and MLOnly produce \emph{identical total FPR} at every seed. At seed=42 MLOnly has 22
hard-blocks plus 90 flags ($112/200$, 56.0\%) and ACB-Ref has 0 hard-blocks plus 112 flags
(also 56.0\%): the rules-certainty veto converts every auto-block into a reviewable flag
without changing the review queue, since LocalRules scoring $0$ below the \$3,000 ceiling is
positive evidence of no deterministic violation. MLOnly's 11.0\% hard FPR comes from the
log-amount overlap between high-value travel (\$1,800--\$2,800) and marginal B8 overspend
(\$3,010--\$4,500); RulesOnly avoids it entirely but pays 58.8\% FNR, since B2/B4/B5/B6 carry
no signal a rule can read. So ACB-Ref's F1 gain over MLOnly is purely recall, and its
operational benefit is the \textbf{elimination of auto-blocks} rather than less review. The
56\% EscRate is the price: those sessions cannot be auto-approved without losing B8 recall,
though a reviewer clears a soft flag in seconds against the irreversibility of a missed
overspend, and per-persona thresholds should reduce it. Errors are at least well-signposted
--- all 37 missed frauds fall in the mid-uncertainty bucket and all 112 false positives in
the high-uncertainty bucket, so hard cases are marked uncertain rather than confidently
mislabeled.

\subsection{External Validation and Open Gaps}
\label{sec:cross_bench}
\label{sec:gaps}

Two external checks guard against the benchmark being self-confirming. On \textbf{AgentDojo},
we ran ACB-Ref's B6 scanner unchanged over the eight canonical Banking injection templates
(\texttt{important\_instructions})~\cite{agentdojo}: 8/8 detected, 0/10 false positives on
clean metadata, so the injection class transfers to templates we did not write. On
\textbf{$\tau$-bench}, mapping the 115 retail trajectories~\cite{taubench} onto ACB action
types, 92.2\% violate the strict automaton $\mathcal{M}$ --- scope mismatch rather than
over-flagging, since $\tau$-bench covers post-purchase servicing (exchange, return, cancel)
rather than purchase flows. Extending to $\mathcal{M}'$ with FIND$\to$COMMIT for service
operations drops this to 13.9\%, locating ACB's boundary precisely: it models purchase-side
commerce, not post-purchase servicing.

The gaps we do not close are as follows. \textbf{B4 (68\%)} needs a persistent cross-session
agent-ID registry and \textbf{B5 (70\%)} a stateful settled-key registry; we rely on
correlated ML signal for both. \textbf{B7s (82.5\%)} leaves MCC--merchant mismatch inside the
model's uncertainty band, needing merchant-domain knowledge unavailable at L1.
\textbf{EscRate (56\%)} needs per-persona amount thresholds to reduce without costing B8
recall. \textbf{B9 (revert-grant)} is defined but excluded, requiring post-settlement rail
state. Finally, B7s/B8s show behavioral variants of rule-trivial attacks shift the problem
substantially (rules 0\%, ACB-Ref 82.5\%); a planned v2 adds MCC aliasing, near-valid grammar
sequences, and boundary-straddle sessions so that no method saturates any class.

% ─────────────────────────────────────────────────────────────────────────────
\section{Conclusion}
\label{sec:conclusion}

Structural violations fall to deterministic rules at 0\% FPR and need no benchmark. ACB
targets the hard part---behavioral fraud that is invisible to reasoning-layer tools and
genuinely ambiguous against legitimate high-value traffic---and scores both defensive layers
on identical sessions. The central finding is that they are complementary rather than
ordered: L0 reaches 61.0\% and ACB-Ref 89.0\%, both at 0\% hard FPR, yet L0 strictly wins on
B6 and B7s while L1 wins on six classes, and the per-attack maximum of 93.3\% exceeds either.
Instrumenting one layer leaves attacks the other would have caught, and no deployed system
spans both. Generator, verifiers, and harness are released so the gaps above can be attacked
directly.

\begin{thebibliography}{99}
\small

\bibitem{garak} Derczynski, L., Galinkin, E., Martin, J., Majumdar, S., \& Inie, N.
``garak: A Framework for Security Probing Large Language Models.'' \emph{arXiv:2406.11036}, 2024.

\bibitem{promptfoo} promptfoo contributors. ``promptfoo: Open-source LLM evaluation and
red-teaming framework.'' \url{https://github.com/promptfoo/promptfoo}, accessed 2026-09-06.

\bibitem{llamaguard} Inan, H., Upasani, K., Chi, J., Rungta, R., Iyer, K., Mao, Y.,
Tontchev, M., Hu, Q., Fuller, B., Testuggine, D., \& Khabsa, M. ``Llama Guard: LLM-based
Input-Output Safeguard for Human-AI Conversations.'' \emph{arXiv:2312.06674}, 2023.

\bibitem{finharness} Jia, H., Liu, Y., Chong, B., Yang, Y., Chen, Y., Liang, J., Li, Q.,
Lu, H., Xu, K., Zheng, H., Zhang, C., Peng, H., \& Yu, P.\ S. ``FinHarness: An Inline
Lifecycle Safety Harness for Finance LLM Agents.'' \emph{arXiv:2605.27333}, 2026.

\bibitem{x402} Reppel, E., Caspers, R., Leffew, K., Organ, D., Kim, D., \& Dalal, N.
``x402: An Open Standard for Internet-Native Payments.'' Coinbase Developer Platform
whitepaper, 2025. \url{https://www.x402.org/x402-whitepaper.pdf}

\bibitem{x402attacks} Li, Z., Wang, Q., \& Wang, Z. ``Five Attacks on x402 Agentic Payment
Protocol.'' \emph{arXiv:2605.11781}, 2026.

\bibitem{agentdojo} Debenedetti, E., Zhang, J., Balunovi\'{c}, M., Beurer-Kellner, L.,
Fischer, M., \& Tram\`{e}r, F. ``AgentDojo: A Dynamic Environment to Evaluate Prompt
Injection Attacks and Defenses for LLM Agents.'' \emph{Advances in Neural Information
Processing Systems 37 (NeurIPS 2024), Datasets and Benchmarks Track}; \emph{arXiv:2406.13352}.

\bibitem{iso8583} ISO 8583:2023. \emph{Financial-transaction-card-originated messages ---
Interchange message specifications}, 3rd ed. International Organization for Standardization, 2023.

\bibitem{stripe} Stripe Inc. ``Payment Intents API reference.''
\url{https://docs.stripe.com/api/payment_intents}, accessed 2026-09-06.

\bibitem{taubench} Yao, S., Shinn, N., Razavi, P., \& Narasimhan, K. ``$\tau$-bench: A
Benchmark for Tool-Agent-User Interaction in Real-World Domains.'' \emph{arXiv:2406.12045}, 2024.

\bibitem{agentharm} Andriushchenko, M., Souly, A., Dziemian, M., Duenas, D., Lin, M.,
Wang, J., Hendrycks, D., Zou, A., Kolter, Z., Fredrikson, M., Winsor, E., Wynne, J.,
Gal, Y., \& Davies, X. ``AgentHarm: A Benchmark for Measuring Harmfulness of LLM Agents.''
\emph{ICLR 2025}; \emph{arXiv:2410.09024}.

\bibitem{pozzolo2015} Dal Pozzolo, A., Caelen, O., Johnson, R.\ A., \& Bontempi, G.
``Calibrating Probability with Undersampling for Unbalanced Classification.''
\emph{IEEE Symposium Series on Computational Intelligence (SSCI)}, 2015, pp.~159--166.

\end{thebibliography}

\begin{ack}
This work was conducted without external funding. Competing interests: none.
\end{ack}

\section*{Reproducibility Statement}

All results are deterministic given \texttt{--seed}. Generator
(\texttt{benchmark/synthetic.py}, including B7s/B8s behavioral variants),
verifiers (\texttt{benchmark/acp\_baselines.py}, \texttt{benchmark/behavioral\_ml.py}),
and harness (\texttt{benchmark/run\_experiment.py}) are released as open-source.
Replication: \texttt{python -m benchmark.run\_experiment --seed 42 --no-sota}. No GPU required;
$<$30s per seed on CPU. Results saved to \texttt{benchmark/results/acp\_benchmark.json}.
Full-session LLM-as-judge replication: \texttt{python -m benchmark.session\_fraud\_judge
--all-attacks --seed 42 --n 20} (requires AWS Bedrock; $\approx$7 min, \$0.60).

\section*{Broader Impact}

ACB targets the protocol layer between LLM reasoning and payment rails.
Benefits: better-defended agentic payment pipelines; a shared benchmark for
protocol-layer safety research; a vocabulary for previously underdescribed
threats. Risks: the attack taxonomy could inform adversarial probing; mitigated
by releasing the verifier alongside the taxonomy and withholding working exploit
code. The multi-class structure (ten orthogonal attacks spanning rule-trivial,
behavioral, stateful, and distributional challenge classes) makes benchmark
gaming substantially harder than single-surface attacks.


\appendix

\section{Benchmark Generator Details}

The full generator is available at \url{https://github.com/[anonymous]/acb}. Key implementation notes:

\paragraph{B5 idempotency key format.} Attack keys use the format \texttt{sess-clean-NNNN:act:2} where \texttt{NNNN} is a prior settled session index. This is format-identical to legitimate keys---the only detection mechanism is a stateful settled-key registry, which the current rules verifier does not maintain cross-session.

\paragraph{B7 MCC categories.} Six varied restricted categories: gambling (7995), pharmacy (5912), liquor store (5921), matchmaking app (7273), wellness spa (7297), currency exchange (6051). Using a single hardcoded merchant would make the attack trivially detectable by keyword match.

\paragraph{B8 bimodal distribution.} 60\% of B8 sessions sample from $\text{Uniform}(\$3{,}010, \$4{,}500)$---marginal overspend above the \$3,000 ceiling. 40\% from $\text{Uniform}(\$4{,}500, \$8{,}000)$---obvious overspend. The marginal fraction creates genuine overlap with high-value clean sessions.

\paragraph{Reproducibility.} Run: \texttt{python -m benchmark.run\_experiment --seed 42 --no-sota}

\section*{NeurIPS Paper Checklist}

\begin{enumerate}

\item {\bf Claims}
    \item[] Question: Do the main claims made in the abstract and introduction accurately reflect the paper's contributions and scope?
    \item[] Answer: \answerYes{}
    \item[] Justification: The abstract reports multi-seed mean $\pm$ std results (seeds 42--46) that match Table~\ref{tab:full} exactly. The per-attack Table~\ref{tab:detection} clearly distinguishes multi-seed rows (ML only, ACB-Ref; seeds 42--46, 520 sessions, 8 attacks) from seed=42 supplementary rows (B7s, B8s, full-session LLM judge; labeled in caption). The contributions list in \S\ref{sec:intro} matches the content of \S\ref{sec:related}--\S\ref{sec:results}. The full-session LLM-as-judge results (61\% recall, 0\% hard FPR) are explicitly labeled as measured on ACB at seed=42. The total-FP equivalence finding is disclosed in the abstract and \S\ref{sec:results}.

\item {\bf Limitations}
    \item[] Question: Does the paper discuss the limitations of the work performed by the authors?
    \item[] Answer: \answerYes{}
    \item[] Justification: \S\ref{sec:gaps} (Open Gaps) covers: (1) B4 at 68\% requires an external agent-ID registry; (2) B5 at 70\% uses ML behavioral signal rather than a stateful settled-key registry; (3) B7s at 82.5\% requires merchant-domain knowledge beyond the current MCC allowlist; (4) EscRate (56\%) is a per-persona spend-ceiling design issue, not a calibration artifact---ACB-Ref hard FPR is exactly 0\%; (5) B9 excluded from this release; (6) B7s/B8s results are seed=42 only, not yet multi-seed. The full-session LLM judge (61\% recall) is described as an L0 ceiling subject to observation-surface limits, not as a deployable component.

\item {\bf Theory assumptions and proofs}
    \item[] Question: For each theoretical result, does the paper provide the full set of assumptions and a complete proof?
    \item[] Answer: \answerNA{}
    \item[] Justification: This is primarily a benchmark and systems paper. Definition~2.1 is a formal definition, not a theorem. No theorems are stated or proved. The paper's claims are empirical and definitional; any formal development is out of scope for this submission.

\item {\bf Experimental result reproducibility}
    \item[] Question: Does the paper fully disclose all information needed to reproduce the main experimental results?
    \item[] Answer: \answerYes{}
    \item[] Justification: The Reproducibility Statement gives the exact reproduction command for the core benchmark (\texttt{python -m benchmark.run\_experiment --seed 42 --no-sota}; CPU only, $<$30s/seed) and for the full-session LLM judge (\texttt{python -m benchmark.session\_fraud\_judge --all-attacks --seed 42 --n 20}; requires AWS Bedrock, $\approx$7 min, \$0.60). Markov parameters, amount distributions, perturbation $\sigma$, training seed, threshold formula ($\mu + 1.1\sigma$), and veto conditions are specified in \S\ref{sec:method}, \S\ref{sec:gen}, and the appendix. B7s/B8s attack generator is in \texttt{benchmark/synthetic.py}.

\item {\bf Open access to data and code}
    \item[] Question: Does the paper provide open access to data and code?
    \item[] Answer: \answerYes{}
    \item[] Justification: Generator, verifiers, and harness are released as open-source (anonymized URL in supplemental material; full URL in camera-ready). The benchmark is fully deterministic given \texttt{--seed}; no pre-collected dataset is required.

\item {\bf Experimental setting/details}
    \item[] Question: Does the paper specify all training and test details necessary to understand the results?
    \item[] Answer: \answerYes{}
    \item[] Justification: Data generation: persona-conditional Markov chains with parameters in \S\ref{sec:gen} and the appendix; perturbation $\sigma=0.05$. Train/test independence: ML trained on seed=7 ($N=500$), evaluated on seeds 42--46 ($N=520$ each). Threshold: $\mu+1.1\sigma$ of clean training scores. Veto: score$=0$ and uncertainty$>0.20$. All in \S\ref{sec:method} and \S\ref{sec:experiments}.

\item {\bf Experiment statistical significance}
    \item[] Question: Does the paper report error bars or appropriate information about statistical significance?
    \item[] Answer: \answerYes{}
    \item[] Justification: Table~\ref{tab:full} reports mean $\pm$ standard deviation across 5 independent seeds (42--46). The $\pm$ values are population standard deviations, not standard errors of the mean---stated explicitly in the table caption. All five seeds are fully independent (separate data generation, no shared state). Rates are bounded well away from 0 and 1 in the reported range, so symmetric error bars do not produce out-of-range values.

\item {\bf Experiments compute resources}
    \item[] Question: Does the paper provide sufficient information on compute resources needed to reproduce the experiments?
    \item[] Answer: \answerYes{}
    \item[] Justification: Stated in the Reproducibility Statement: CPU only (no GPU required), $<$30 seconds per seed, $<$3 minutes total for all 5 seeds. The Tier~3 LLM Verifier is invoked on $<$3\% of actions in the full pipeline but is bypassed by \texttt{--no-sota} for the main results; \texttt{--no-sota} results are the ones reported in all tables.

\item {\bf Code of ethics}
    \item[] Question: Does the research conform with the NeurIPS Code of Ethics?
    \item[] Answer: \answerYes{}
    \item[] Justification: This work promotes defensive safety in financial AI systems. All data is synthetically generated---no production transaction data or user data was used. No production payment credentials appear in the open-source release. The attack taxonomy is published at a conceptual level sufficient to motivate defenses, not at a working-exploit level.

\item {\bf Broader impacts}
    \item[] Question: Does the paper discuss both potential positive and negative societal impacts?
    \item[] Answer: \answerYes{}
    \item[] Justification: See the Broader Impact section following the Conclusion. Positive: protocol-layer safety tools may reduce financial fraud in agentic systems. Negative: the attack taxonomy, if studied adversarially, could inform attack design; mitigated by not publishing working exploit code and by releasing the detector alongside the taxonomy.

\item {\bf Safeguards}
    \item[] Question: Does the paper describe safeguards for responsible release of data or models with high risk for misuse?
    \item[] Answer: \answerYes{}
    \item[] Justification: ACB describes attack \emph{classes} at a conceptual and statistical level (Markov chains, amount distributions). No production payment API credentials, no merchant exploit scripts, and no working B5 settled-key replay against a live rail are included in the release. The released artifact is the detector (ACB-Ref), not the attack implementation. The benchmark generator produces synthetic sessions only.

\item {\bf Licenses for existing assets}
    \item[] Question: Are creators of assets used in the paper properly credited and license terms mentioned?
    \item[] Answer: \answerYes{}
    \item[] Justification: All assets cited are referenced, not redistributed --- no third-party code is reused in ACB. garak~\cite{garak} (Apache 2.0) and promptfoo~\cite{promptfoo} (MIT) are named as related tooling and are not run or bundled; Llama Guard~\cite{llamaguard} and FinHarness~\cite{finharness} are cited as prior work only. ISO~8583~\cite{iso8583} is a published standard (referenced, not reproduced); x402~\cite{x402}, its security analysis~\cite{x402attacks}, and the Stripe API reference~\cite{stripe} are cited as public specifications and documentation. The benchmarks we cross-validate against, AgentDojo~\cite{agentdojo} and $\tau$-bench~\cite{taubench}, are used under their published open-source licenses (MIT), with attack templates and task trajectories credited in \S\ref{sec:cross_bench}. All citations appear in the References section and were verified against their primary sources.

\item {\bf New assets}
    \item[] Question: Are new assets introduced in the paper well documented?
    \item[] Answer: \answerYes{}
    \item[] Justification: ACB (benchmark generator, verifiers, harness) is documented in \S\ref{sec:dataset} (methodology), the appendix (implementation details), and the Reproducibility Statement (exact run command). Released under MIT license. Anonymous supplemental URL at submission; full URL in camera-ready.

\item {\bf Crowdsourcing and research with human subjects}
    \item[] Question: Does the paper include full details of crowdsourcing or human subjects research?
    \item[] Answer: \answerNA{}
    \item[] Justification: All data is synthetically generated via parameterized Markov chains and log-normal amount distributions. No human subjects, no crowdsourcing, no real user data.

\item {\bf Institutional review board (IRB) approvals}
    \item[] Question: Does the paper describe potential risks to study participants and IRB approvals?
    \item[] Answer: \answerNA{}
    \item[] Justification: No human subjects research. Not applicable.

\item {\bf Declaration of LLM usage}
    \item[] Question: Does the paper describe the usage of LLMs if it is an important, original, or non-standard component of the core methods?
    \item[] Answer: \answerYes{}
    \item[] Justification: Two LLM uses are disclosed. (1) The Tier~3 LLM Verifier (\S\ref{sec:method}) uses \texttt{claude-haiku-4-5-20251001} as a classification component for high-uncertainty actions ($<3\%$ of actions in the full pipeline)---a method component described in \S\ref{sec:method}. (2) The full-session LLM-as-judge baseline uses \texttt{claude-sonnet-4-6} via AWS Bedrock to evaluate all 10 attack types ($n{=}20$/attack, seed=42, reported in Table~\ref{tab:detection})---an evaluation baseline, not a method component. LLMs are \emph{not} used in benchmark generation, L1 rules verification, or the ML behavioral model. LLMs were also used for writing assistance in preparing this manuscript, which does not require declaration per the NeurIPS handbook.

\end{enumerate}


\end{document}
